% ============================================================================
%  lessons/lesson03.tex — قالبِ آمادهٔ درجِ محتوا (درسِ سوم)
% ============================================================================
\Lesson{اَلدَّرْسُ الثَّالِثُ}{ثَلَاثُ قِصَصٍ قَصِيرَةٍ}{[حکمتِ عربیِ آغازِ درس]}{[ترجمهٔ حکمت]}{[نقل‌کننده]}
\JGoals{هدف‌های درسِ سوم اینجا درج شود — مثلاً: حکایت‌خوانی و ترجمه \JGoalSep قواعد درس}

% ------------------------------------------------------- متنِ درس -----------
\Jsec{اَلنَّصُّ \textbar\ متنِ درس}
\begin{JNass}
(متنِ درسِ سوم اینجا درج شود: عنوانِ قصهٔ نخست، ابیات یا بندهای نثر،
با پانویس‌های خودِ کتاب. برای شعر از محیطِ \texttt{JPoem} و برای هر بیت
از دستورِ \texttt{JBayt} استفاده کنید.]
\end{JNass}

% ------------------------------------------------------- ترجمهٔ درس ---------
\Jsec{اَلتَّرْجَمَةُ \textbar\ ترجمهٔ روان}
\begin{JTarjome}
(ترجمهٔ روانِ متنِ درسِ سوم اینجا درج شود.)
\end{JTarjome}

% ------------------------------------------------------- واژه‌نامه ----------
\Jsec{اَلْمُعْجَمُ \textbar\ واژه‌نامهٔ درس}
\begin{JMTable}
\JW{اَلْكَلِمَةُ}{اَلْمَعْنَى}{روابط/نکته}
\end{JMTable}

% ------------------------------------------------------- قواعد درس ----------
\Jsec{إِعْلَمُوا \textbar\ آموزشِ قواعد}
\begin{JQaede}[JCompact]{عنوانِ قاعدهٔ درسِ سوم}
(شرحِ کاملِ قاعده، جدول‌ها و مثال‌ها اینجا درج شود.]
\end{JQaede}

% ------------------------------------------------------- خودآزمایی ----------
\Jsec{اِخْتِبَارْ نَفْسَكَ \textbar\ خودآزمایی}
\begin{JTest}[خودآزماییِ ۱]
(پرسش‌های خودآزماییِ کتاب اینجا درج شود.)
\end{JTest}

% ------------------------------------------------------- تمارین کتاب --------
\Jsec{اَلتَّمَارِينُ \textbar\ تمرین‌های کتاب}
\begin{JTamrin}{۱}{عنوانِ تمرین}
(متنِ تمرینِ نخست اینجا درج شود.)
\end{JTamrin}

% ------------------------------------------------------- پاسخ‌نامه ----------
\Jsec{اَلْإِجَابَاتُ \textbar\ پاسخ‌نامهٔ خودآزمایی و تمرین‌ها}
\begin{JKey}[پاسخ‌نامهٔ خودآزمایی]
\JAnswer{۱}{(پاسخِ تشریحی اینجا درج شود.)}
\end{JKey}
\begin{JKey}[پاسخ‌نامهٔ تمرین‌ها]
\JAnswer{تمرینِ ۱}{(پاسخ اینجا درج شود.)}
\end{JKey}
